\documentclass{minmal}
\usepackage{array}
\renewcommand{\arraystretch}{1.5}
\begin{document}
\begin{center}
    \textbf{Løsningsforslag til prøve i mekaniske separasjonsmetoder:}
\end{center}
\noindent
\textbf{Oppg. 1:}\\\\
\noindent
a) Hva menes en homogen og en heterogen blanding? Gi et eksempel for hver av dem.\nl[2]
\textbf{\textit{\textbf{\textit{Homogene blandinger:}}}}

\noindent – er blandinger hvor atomene eller molekylene stoffene består av er likt fordelt. Blandingen er så jevn at vi ikke kan skille eller se hvilke stoffer som er i blandingen. Eksempler på homogene blandinger er:

\begin{itemize}
    \item Saltløsninger – som i sjøvann;
    \item Metall legeringer;
    \item Gassblandinger – som i luft.\\\\
\end{itemize}
\textbf{\textit{\textit{Heterogene blandinger:}}}

– er blandinger som består av partikler eller partikkelstørrelser, som korn eller dråper, som er fint fordelt i blandingen. \textit{Heterogene blandinger består av ulike stoffer som egentlig ikke er blandbare med hverandre}, og du kan ane at blandingen består av ulike stoffer. Stoffene som er blandet kan ha ulike fasetilstander. Eksempler på heterogene blandinger er:

\begin{itemize}
    \item Pulverblandinger (tørrstoff blandinger);
    \item Suspensjoner (deiger og slurry);
    \item Emulsjoner (pasta og kremer);
    \item Aerosol (røyk).
\end{itemize}
 \newpage
\noindent b) Hva menes med en suspensjon? Gi et eksempel på en suspensjon. \nl[1]\\
\textbf{Svar:} En suspensjon er en faststoff-væ'skeblanding. Eksmepel kan være avløpsvann og papirmasse.\nl[2]

\noindent c) Hva vil skje med partikler i vann som har større tetthet enn vann, og hva vil skje med partikler i vann som har tetthet mindre enn vann?\nl[1]
\textbf{Svar:} For partikler i en væske så er det slik at hivs partiklene har større massetetthet enn væsken (i dette tilfellet vann) så vil de synke. Partikler med massetetthet mindre enn væsken (i dette tilfellet vann) vil flyte opp. \\\\
d) Hvilke to emulsjoner har vi?  \nl[1]
\textbf{Svar:} Vi har to typer emulsjoner, og det er vann-i-olje og olje-i vann.\\\\
e) Hva menes med mettet, umettet og overmettet løsning. Bruk $NaCl$ som eksempel.  Ved $25^0C$ for du maksimum løst 360g/L. \nl[1]
\textbf{Svar:} Ved $25^0C$ kan man få løst opp 360g $NaCl$ i 1L vann. Når vi har 360g så sier vi at løsningen er mettet.  Har vi mindre enn 360g så sier vi at løsningen er umettet.  Og er vi over 360g så sier vi at løsningen er overmettet.  Da vil mengde salt over 360g foreligge som fast stoff, bunnfall. \newpage

\noindent \textbf{Oppg. 2:}\\\\
a) Figuren under viser en standard rørebeholder:\\

\begin{tabular}{m{8cm}>{\raggedright\arraybackslash}m{6cm}}

\includegraphics[width=8cm]{rorebeholder.png} &
1) Hva brukes inspeksjonsluken til? \newline \newline
2) Hva er oppgaven til temperaturføleren? \newline \newline
3) Hva er oppgaven til elektromotoren? \\

\end{tabular}\\\noindent \textbf{Svar:} \\
\textbf{Inspeksjonsluken} brukes til å inspisere tanken innvendig. Man kan også veldig ofte gå inn i tanken via denne luken for å gjøre vedlikeholdsarbeid. \\\\
\textbf{Temperaturføler:} Brukes til å måle temperaturen til væsken.\\\\
\textbf{Elektromotoren:} Sørger for å få røreren til å gå rundt.\nl[2]
b) Forklar kort hva som menes med begrepet viskositet for en væske. Gi et eksempel på en væske med høy viskositet og en med lav viskositet.\nl[1]
\textbf{Svar:} En væskes viskositet angir hvor seig en væske er. Har væsken høy viskositet er den tyktflytene (f.eks. sirup).  Er væsken tyntflytende har den lav viskositet (f.eks.  vann). \newpage

c) Det finnes mange typer rørere:
\begin{figure}[h]
    \centering
    \includegraphics[width=0.9\linewidth]{rørere.png}
    \caption{Ulike rørere}
    \label{fig:placeholder}
\end{figure}
En propellrører er hurtiggående, en spiral Wendel er saktegående. hvilken av de to rørerne vil du bruke i en sirup-tykk løsning? Begrunn svaret.\nl[2]
\textbf{Svar:} Når man skal røre i tygtflytende væske (en væske med høy viskositet) så må man ha en saktegående rører. Derfor må man velge en spiral Wendel rører.  
\textbf{Oppg. 3:}\nl[1]
a) Forklar forskjellen på sortering og klassering i kjemisk teknologi.\nl[2]
\textbf{Svar:}\\\\
\textbf{Sortering} handler om å skille to (eller flere) stoffer i en blanding etter type stoff. \\\\
\textbf{Klassering} handler om å separere stoffer etter størresle.\newpage

b) Forklar virkemåten til en hydrosyklon:
\begin{figure}[h]
    \centering
    \includegraphics[width=0.65\linewidth]{hydrosyklon.png}
    \caption{En hydrosyklon}
    \label{fig:placeholder}
\end{figure}\\
En hydrosyklon brukes til separering av suspensjoner som inneholder faststoffpartikler med ulik tetthet.  Inne i syklonen vil partiklene og væsken bevege seg i en sirkulær bane slik at det virker en sentrifugalkraft på partiklene. Hastigheten må være slik at partiklene med størst tetthet treffer veggen og sklir ned mot bunnen hvor de tas ut.  Partiklene med lavest tetthet og vannet vil tilslutt gå i en indre spiral og forlate syklonen på toppen.\newpage
c) Forklar virkemåten til magnetsortereren:\\
\begin{figure}[h]
    \centering
    \includegraphics[width=0.7\linewidth]{magnetseparator.png}
    \caption{magnetsorterer}
    \label{fig:placeholder}
\end{figure}\\
\textbf{Svar :}Stoffet som føres til magnetsortereren består av partikler som er magnetiske, og partikler som ikke er magnetiske.  De føres over et vibrasjonsbelte og så mot en stor magnet.  Stoffet som ikke er magnetisk farer rett frem og ned i container.  Stoffet som er magnetisk vil følge med videre og føres til en annen container. \newpage

\noindent d) Forklar virkemåten til et rundt klaringsbasseng:
\begin{figure}[h]
    \centering
    \includegraphics[width=0.84\linewidth]{rundt_klaringsbasseng.png}
    \caption{Et rundt klaringsbasseng}
    \label{fig:placeholder}
\end{figure}
\textbf{Svar:} et rundt klaringsbasseng er ofte stort. Det kan ha en diameter på ca. 20m og en dybde på ca. 10m.  Her tilføres det en suspensjon med partikler som ALLE har tetthet større enn vann. I dette bassenget har man derfor IKKE flotasjon.   Det rensede Vannet vil forlate beholderen i overløp\\\\
e) Forklar virkemåten til flotasjonsfortykkeren:
\begin{figure}[h]
    \centering
    \includegraphics[width=0.9\linewidth]{flotasjonsfortykker.png}
    \caption{Flotasjonsfortykker}
    \label{fig:placeholder}
\end{figure}\nl[1]
\textbf{Svar:} Flotasjonsfortykkeren brukes til rensing av avløpsvann.  suspensjonen består av faststoff med tetthet større enn vann og partikler med tetthet mindre enn vann.  I tanken blir suspensjon (væske med faststoff) ført inn i toppen av tanken. Suspensjonen blir ledet inn i et «skjørt», slik at væske får retning nedover. Det gir raskere sedimentering av stoffer med høyest tetthet. Med litt tid vil faststoff med tetthet større enn vann bli tatt ut som slam i bunnen av tanken. Faststoffpartiklene med tetthet mindre enn vann vil flyte med væsken opp og legge seg som et lag over væsken. Dette blir tatt ut i et overløp eller skrapet av.  Det er viktig å ha en viss høyde på denne for å få skilt stoffene. \\\\
f) For rensing av gass bruker man et støvkammer:
\begin{figure}[h]
    \centering
    \includegraphics[width=0.8\linewidth]{stovkammer.png}
    \caption{Støvkammer}
    \label{fig:placeholder}
\end{figure}\nl[1]
Forklar virkemåten til støvkammeret.\\\\
\textbf{Svar:}  Rågass med støvpartikler føres inn i støvkammeret.  Siden støvkammeret er stort bremses gassen og støvpartiklene. De vil derfor falle mot bunnen av støvkammeret og tas ut der.  Det er også plater inne i støvkammeret. Når støvpartiklene treffer disse vil de falle ned.
\end{document}